%%%%%%%%%%%%%%%%%%%%%%%%%%%%%%%%%%%%%%%%%%%%%%%%%%%%%%%%%%%%%%%%%%%%%%%%%%%%%%
%% Preambolo -- Algoritmi Avanzati
%% Classe: academica (Airy).  Compilare con lualatex.
%%%%%%%%%%%%%%%%%%%%%%%%%%%%%%%%%%%%%%%%%%%%%%%%%%%%%%%%%%%%%%%%%%%%%%%%%%%%%%

% academica carica gia': tikz, pgfplots, mathtools, amssymb, tcolorbox, xcolor,
% graphicx, float, caption/subcaption, listings, hyperref, xspace, unicode-math.

\usetikzlibrary{automata,positioning,trees,arrows.meta,calc,fit,shapes.geometric,
                decorations.pathreplacing,decorations.markings,matrix,backgrounds,chains}
\usepackage{booktabs}
\usepackage{multirow}
\usepackage{array}
\usepackage{algorithm}
\usepackage{algpseudocode}
\usepackage{marginnote}

% Localizzazione italiana.
% NB: non si puo' usare babel/polyglossia perche' academica carica biblatex
% (che pretende di essere caricato DOPO il pacchetto di lingua). Rinominiamo
% quindi a mano le poche stringhe che compaiono nel documento.
\renewcommand{\chaptername}{Capitolo}
\renewcommand{\contentsname}{Indice}
\renewcommand{\listfigurename}{Elenco delle figure}
\renewcommand{\listtablename}{Elenco delle tabelle}
\renewcommand{\figurename}{Figura}
\renewcommand{\tablename}{Tabella}
\renewcommand{\bibname}{Bibliografia}
\renewcommand{\indexname}{Indice analitico}
\renewcommand{\appendixname}{Appendice}
\renewcommand{\partname}{Parte}

\providecommand{\qedhere}{}

% academica scrive "List of Figures"/"List of Tables" dentro una tikzpicture,
% quindi \listfigurename non basta: va ridefinito il comando.
\makeatletter
\renewcommand{\listoffigures}{%
  \chapter*{%
    \vspace*{-20\p@}%
    \begin{tikzpicture}[overlay]
      \draw[fill=accent,draw=accent,rounded corners] (10,-0.65) rectangle (20,1);
      \pgftext[right,x=16.7cm,y=0.2cm]{\Huge\sc\bfseries\textcolor{white}{\listfigurename}}
    \end{tikzpicture}%
  }%
  \@starttoc{lof}%
}
\renewcommand{\listoftables}{%
  \chapter*{%
    \vspace*{-20\p@}%
    \begin{tikzpicture}[overlay]
      \draw[fill=accent,draw=accent,rounded corners] (10.5,-0.65) rectangle (20,1);
      \pgftext[right,x=16.7cm,y=0.2cm]{\Huge\sc\bfseries\textcolor{white}{\listtablename}}
    \end{tikzpicture}%
  }%
  \@starttoc{lot}%
}
\makeatother

\setlength{\parindent}{0pt}

%%%%%%%%%%%%%%%%%%%%%%%%%%%%%%%%%%%%%%%%%%%%%%%%%%%%%%%%%%%%%%%%%%%%%%%%%%%%%%
%% AMBIENTI IN ITALIANO
%% academica definisce gli ambienti colorati in inglese (Definition, Theorem,
%% ...). Qui aggiungiamo gli alias italiani, numerati e non, riusando gli
%% stessi contatori e colori della classe.
%%%%%%%%%%%%%%%%%%%%%%%%%%%%%%%%%%%%%%%%%%%%%%%%%%%%%%%%%%%%%%%%%%%%%%%%%%%%%%

\makeatletter
% versione numerata (riusa i contatori della classe: stessa numerazione
% degli ambienti inglesi, così i due stili non divergono)
\newcommand{\aa@num}[4]{% #1 nome env, #2 contatore, #3 etichetta, #4 colore
  \newenvironment{#1}[1][]{%
    \refstepcounter{#2}\def\@currentlabelname{##1}%
    \begin{accentbox}[##1]{#3}{\csname the#2\endcsname}{#4}%
  }{\end{accentbox}}%
  \expandafter\newcommand\csname #1autorefname\endcsname{#3}%
}
\aa@num{definizione}{definition}{Definizione}{lilla}
\aa@num{teorema}{theorem}{Teorema}{blue}
\aa@num{proposizione}{proposition}{Proposizione}{deepsea}
\aa@num{corollario}{corollary}{Corollario}{deepsea}
\aa@num{dimostrazione}{proof}{Dimostrazione}{pink}
\aa@num{esempio}{example}{Esempio}{orange}
\aa@num{esercizio}{exercise}{Esercizio}{orange}
\aa@num{domanda}{question}{Domanda}{green}
\makeatother

% versioni non numerate
\newenvironment{definizione*}[1][]{\begin{accentbox*}[#1]{Definizione}{lilla}}{\end{accentbox*}}
\newenvironment{teorema*}[1][]{\begin{accentbox*}[#1]{Teorema}{blue}}{\end{accentbox*}}
\newenvironment{proposizione*}[1][]{\begin{accentbox*}[#1]{Proposizione}{deepsea}}{\end{accentbox*}}
\newenvironment{corollario*}[1][]{\begin{accentbox*}[#1]{Corollario}{deepsea}}{\end{accentbox*}}
\newenvironment{dimostrazione*}[1][]{\begin{accentbox*}[#1]{Dimostrazione}{pink}}{\end{accentbox*}}
\newenvironment{esempio*}[1][]{\begin{accentbox*}[#1]{Esempio}{orange}}{\end{accentbox*}}
\newenvironment{esercizio*}[1][]{\begin{accentbox*}[#1]{Esercizio}{orange}}{\end{accentbox*}}
\newenvironment{domanda*}[1][]{\begin{accentbox*}[#1]{Domanda}{green}}{\end{accentbox*}}
\newenvironment{lemma*it}[1][]{\begin{accentbox*}[#1]{Lemma}{deepsea}}{\end{accentbox*}}

% "Osservazione" non numerata (equivalente italiano di \begin{note})
\newenvironment{osservazione}[1][]{\begin{accentbox*}[#1]{Osservazione}{gray}}{\end{accentbox*}}

% autoref in italiano per gli ambienti della classe
\renewcommand{\lemmaautorefname}{Lemma}
\renewcommand{\chapterautorefname}{Capitolo}
\renewcommand{\sectionautorefname}{Sezione}

\hyphenation{suf-fis-so suf-fis-si pre-fis-so pre-fis-si strin-ga strin-ghe}
\hyphenation{al-li-nea-men-to ran-do-miz-za-to ran-do-miz-za-ti com-pres-se}

%%%%%%%%%%%%%%%%%%%%%%%%%%%%%%%%%%%%%%%%%%%%%%%%%%%%%%%%%%%%%%%%%%%%%%%%%%%%%%
%% MARCATORI EDITORIALI
%% Usati SOLO quando il dubbio resta irrisolto. Ogni scelta e' comunque
%% registrata in NOTE-EDITORIALI.md.
%%%%%%%%%%%%%%%%%%%%%%%%%%%%%%%%%%%%%%%%%%%%%%%%%%%%%%%%%%%%%%%%%%%%%%%%%%%%%%

% Riferimento alla pagina della scansione originale (appunti.pdf).
% Usa \marginnote e non \marginpar: quest'ultimo e' un float e va perso se
% invocato dentro un tcolorbox (cioe' dentro definizione/teorema/...).
%
% NB: academica non riserva spazio a margine (\marginparwidth vale 0pt), per cui
% senza queste due righe ogni \scan produce un Overfull \hbox di ~20pt. Il
% margine destro utile e' 44.8pt (597.5 - 72.27 - 2.42 - 478.01), quindi si sta
% larghi con 34+4 = 38pt.
\setlength{\marginparwidth}{34pt}
\setlength{\marginparsep}{4pt}

\newcommand{\scan}[1]{%
  \marginnote{%
    \raggedright\scriptsize\color{gray!140}\textsf{scan\\p.\,#1}%
  }%
}

% Blocco: lettura del manoscritto non determinabile con certezza
\newenvironment{dubbio}[1][Lettura incerta]{%
  \begin{accentbox*}[#1]{Dubbio}{red}%
}{%
  \end{accentbox*}%
}

% Blocco: buco negli appunti, ricostruzione non conclusiva
\newenvironment{lacuna}[1][Lacuna negli appunti]{%
  \begin{accentbox*}[#1]{Lacuna}{yellow}%
}{%
  \end{accentbox*}%
}

% Versioni inline
\newcommand{\dubbioinl}[1]{{\color{red}[?\,\textit{#1}]}}
\newcommand{\lacunainl}[1]{{\color{orange}[\dots\,\textit{#1}]}}

%%%%%%%%%%%%%%%%%%%%%%%%%%%%%%%%%%%%%%%%%%%%%%%%%%%%%%%%%%%%%%%%%%%%%%%%%%%%%%
%% STRINGHE E STRING MATCHING
%%%%%%%%%%%%%%%%%%%%%%%%%%%%%%%%%%%%%%%%%%%%%%%%%%%%%%%%%%%%%%%%%%%%%%%%%%%%%%

\newcommand{\alfa}{\S}                    % alfabeto (\S = \Sigma da academica)
\newcommand{\str}[1]{\texttt{#1}}         % stringa letterale
\newcommand{\sub}[3]{#1[#2\,..\,#3]}      % sottostringa S[i..j]
\newcommand{\suf}[2]{#1[#2\,..\,]}        % suffisso S[i..]
\newcommand{\pre}[2]{#1[1\,..\,#2]}       % prefisso S[1..i]
\newcommand{\conc}{\cdot}                 % concatenazione
\newcommand{\vuota}{\varepsilon}          % stringa vuota

\DeclareMathOperator{\lcp}{lcp}
\DeclareMathOperator{\lcs}{lcs}
\DeclareMathOperator{\lca}{lca}
\DeclareMathOperator{\lcey}{lce}

% KMP / Z
\newcommand{\spv}[1]{\mathit{sp}_{#1}}     % sp_i
\newcommand{\spp}[1]{\mathit{sp}'_{#1}}    % sp'_i
\newcommand{\Zv}[1]{Z_{#1}}                % Z_i
\newcommand{\Zbox}{Z\text{-box}\xspace}
\newcommand{\fail}{\mathit{fail}}

% Boyer-Moore
\newcommand{\bc}[1]{R(#1)}                 % bad character rule
\newcommand{\Lp}[1]{L(#1)}                 % good suffix (strong) rule
\newcommand{\lp}[1]{\ell(#1)}

% Aho-Corasick
\newcommand{\goto}{\mathit{goto}}
\newcommand{\out}{\mathit{output}}
\newcommand{\nv}[1]{n_{#1}}                % failure link

% Suffix tree / trie
\newcommand{\STrie}{\textsf{STrie}\xspace}
\newcommand{\STree}{\textsf{STree}\xspace}
\newcommand{\albero}[1]{\mathcal{T}(#1)}
\newcommand{\slink}{\mathit{sl}}           % suffix link
\newcommand{\etich}[1]{\mathit{label}(#1)}
\newcommand{\cammino}[1]{\mathit{path}(#1)}

% Allineamento / distanze
\newcommand{\dedit}{d_{\mathrm{edit}}}
\newcommand{\dham}{d_{\mathrm{H}}}
\newcommand{\Dm}[2]{D(#1,#2)}
\newcommand{\Vm}[2]{V(#1,#2)}
\newcommand{\gap}{\mathit{gap}}
\newcommand{\indel}{\mathit{indel}}

%%%%%%%%%%%%%%%%%%%%%%%%%%%%%%%%%%%%%%%%%%%%%%%%%%%%%%%%%%%%%%%%%%%%%%%%%%%%%%
%% ALGORITMI RANDOMIZZATI
%%%%%%%%%%%%%%%%%%%%%%%%%%%%%%%%%%%%%%%%%%%%%%%%%%%%%%%%%%%%%%%%%%%%%%%%%%%%%%

\newcommand{\Prob}{\mathbb{P}}
\newcommand{\E}{\mathbb{E}}
\renewcommand{\Pr}{\operatorname{Pr}}
\DeclareMathOperator{\Var}{Var}
\DeclareMathOperator*{\argmin}{arg\,min}
\DeclareMathOperator*{\argmax}{arg\,max}
\newcommand{\indic}[1]{\mathbf{1}_{\{#1\}}}
\newcommand{\RandQS}{\textsc{RandQS}\xspace}
\newcommand{\LasVegas}{Las~Vegas\xspace}
\newcommand{\MonteCarlo}{Monte~Carlo\xspace}
\newcommand{\mincut}{\textsc{MinCut}\xspace}
\newcommand{\contract}{\textsc{Contract}\xspace}
% classi di complessita' randomizzate
\newcommand{\RP}{\mathsf{RP}}
\newcommand{\coRP}{\mathsf{coRP}}
\newcommand{\ZPP}{\mathsf{ZPP}}
\newcommand{\BPP}{\mathsf{BPP}}
\newcommand{\PP}{\mathsf{PP}}
\newcommand{\PSPACE}{\mathsf{PSPACE}}
\newcommand{\NP}{\mathsf{NP}}
\renewcommand{\P}{\mathsf{P}}   % nota: academica definiva \P = \mathcal{P}

%%%%%%%%%%%%%%%%%%%%%%%%%%%%%%%%%%%%%%%%%%%%%%%%%%%%%%%%%%%%%%%%%%%%%%%%%%%%%%
%% STRUTTURE DATI COMPRESSE
%%%%%%%%%%%%%%%%%%%%%%%%%%%%%%%%%%%%%%%%%%%%%%%%%%%%%%%%%%%%%%%%%%%%%%%%%%%%%%

\DeclareMathOperator{\rank}{rank}
\DeclareMathOperator{\select}{select}
\DeclareMathOperator{\access}{access}
\newcommand{\BWT}{\textsf{BWT}\xspace}
\newcommand{\FMindex}{\textsf{FM-index}\xspace}
\newcommand{\SA}{\mathit{SA}}
\newcommand{\LF}{\mathit{LF}}
\newcommand{\entropia}[1]{H_{#1}}

%%%%%%%%%%%%%%%%%%%%%%%%%%%%%%%%%%%%%%%%%%%%%%%%%%%%%%%%%%%%%%%%%%%%%%%%%%%%%%
%% VARIE
%%%%%%%%%%%%%%%%%%%%%%%%%%%%%%%%%%%%%%%%%%%%%%%%%%%%%%%%%%%%%%%%%%%%%%%%%%%%%%

\newcommand{\Oh}{\mathcal{O}}
\newcommand{\Th}{\Theta}
\newcommand{\Om}{\Omega}
\newcommand{\defeq}{\vcentcolon=}
\renewcommand{\emptyset}{\varnothing}

% Ambienti algoritmo in italiano
\algrenewcommand\algorithmicrequire{\textbf{Input:}}
\algrenewcommand\algorithmicensure{\textbf{Output:}}
\algrenewcommand\algorithmicwhile{\textbf{while}}
\algrenewcommand\algorithmicfor{\textbf{for}}
\algrenewcommand\algorithmicif{\textbf{if}}
\algrenewcommand\algorithmicreturn{\textbf{return}}
\floatname{algorithm}{Algoritmo}
